\begin{proofbox}
\textbf{先分清“定义”与“证明”。} 距离和描述本身是椭圆的定义，无须循环证明。
需要证明的是它与标准方程等价。几何图先解释参数，代数再覆盖所有点。

\textbf{方法一：距离的平方差（优先，少做一次根式移项）。}
设 $F_1(-c,0),F_2(c,0)$，$P(x,y)$，记 $r_1=PF_1,r_2=PF_2$。
距离公式给出
\[
r_1^2-r_2^2=(x+c)^2-(x-c)^2=4cx.
\]
若 $r_1+r_2=2a$，两式相除得 $r_1-r_2=2cx/a$，所以
\[
r_1=a+\frac{cx}{a},\qquad r_2=a-\frac{cx}{a}.
\]
代入 $r_2^2=(x-c)^2+y^2$，逐项消去 $-2cx$，得到
\[
a^2+\frac{c^2x^2}{a^2}=x^2+c^2+y^2,
\qquad
\frac{a^2-c^2}{a^2}x^2+y^2=a^2-c^2.
\]
因 $a^2-c^2=b^2>0$，两边除以 $b^2$ 即为标准方程。

\textbf{逆向验证，排除增根。} 若标准方程成立，$|x|\le a$，并且
\[
(x\pm c)^2+y^2=\left(a\pm\frac{cx}{a}\right)^2.
\]
这里 $a\pm cx/a\ge a-c>0$，故可取正平方根，得到
$r_1=a+cx/a,r_2=a-cx/a$，相加为 $2a$。两个方向均成立。
\end{proofbox}
\clearpage
\begin{proofbox}
\textbf{方法二：直接移项平方（熟悉根式运算时使用）。}
从 $\sqrt{(x+c)^2+y^2}=2a-\sqrt{(x-c)^2+y^2}$ 出发。
在满足原定义的点上，右端就是 $PF_1>0$。平方整理为
\[
a^2-cx=a\sqrt{(x-c)^2+y^2}.
\]
再平方得
\[
(a^2-cx)^2=a^2\bigl((x-c)^2+y^2\bigr),
\]
整理为 $(a^2-c^2)x^2+a^2y^2=a^2(a^2-c^2)$。
逆向仍需用方法一中的正号验证，不能仅凭两次平方就声称等价。

\textbf{几何角度：常数为什么必须大于焦距？}
任意点都有 $PF_1+PF_2\ge F_1F_2$，等号恰在闭线段 $F_1F_2$ 上成立。
这一步只用线段最短性质，直接给出下图的边界分类。
\end{proofbox}
